\documentclass[aspectratio=169]{beamer}
\usetheme{Madrid}
\usecolortheme{seahorse}

\usepackage{graphicx}
\usepackage{booktabs}
\usepackage{amsmath}
\usepackage{amssymb}
\usepackage{tabularx}
\usepackage{array}
\usepackage{ragged2e}
\usepackage{bm}

\title[Chapter 13: Sampling Distributions]{\textbf{Chapter 13: Sampling Distributions}}
\subtitle{School of Mathematical and Computational Sciences//University of Prince Edward Island}
\author{Instructor: Ali Raisolsadat}
\date{Part III: Foundations of Inference}

\setbeamertemplate{navigation symbols}{}
\setbeamertemplate{itemize items}[circle]
\setbeamertemplate{enumerate items}[default]
\graphicspath{{./}{figures/}}

\newcommand{\lectureimage}[2][0.90\linewidth]{%
  \IfFileExists{#2}{%
    \includegraphics[width=#1]{#2}%
  }{%
    \fbox{\parbox[c][0.36\textheight][c]{0.88\linewidth}{%
      \centering\small
      Figure file: \texttt{\detokenize{#2}}\\[0.15cm]
      Generate this figure by knitting the revised Chapter 13 RMarkdown notes.
    }}%
  }%
}

\newcommand{\lectureimagefull}[1]{%
  \IfFileExists{#1}{%
    \includegraphics[width=0.98\linewidth,height=0.78\textheight,keepaspectratio]{#1}%
  }{%
    \fbox{\parbox[c][0.70\textheight][c]{0.92\linewidth}{%
      \centering\small
      Figure file: \texttt{\detokenize{#1}}\\[0.15cm]
      Generate this figure by knitting the revised Chapter 13 RMarkdown notes.
    }}%
  }%
}

\begin{document}

\begin{frame}
  \titlepage
\end{frame}

\begin{frame}{Why This Chapter Matters}
A random sample gives information about a population, but a different random sample usually gives a different answer.

\vspace{0.2cm}
Chapter 13 explains how that natural variation behaves. This idea is the bridge from descriptive statistics to statistical inference.

\vspace{0.2cm}
By the end of the chapter, the main questions should be clear.
\begin{itemize}
  \item What is the difference between a parameter and a statistic
  \item What is a sampling distribution
  \item Why do larger samples give more stable results
  \item When can a Normal model be used for a sample mean or a sample proportion
\end{itemize}
\end{frame}

%\begin{frame}{Chapter Roadmap}
%The material in this chapter follows a natural progression.
%\begin{enumerate}
%  \item Parameters and statistics
%  \item Population distributions, sample distributions and sampling distributions
%  \item The sampling distribution of $\bar{x}$
%  \item The effect of sample size
%  \item The Central Limit Theorem
%  \item The sampling distribution of $\hat{p}$
%  \item The Law of Large Numbers
%  \item Probability examples that prepare for statistical inference
%\end{enumerate}
%\end{frame}

\section{Parameters and Statistics}

\begin{frame}{Parameters and Statistics}
A \textbf{parameter} is a numerical characteristic of a population.

\vspace{0.2cm}
A \textbf{statistic} is a numerical characteristic calculated from sample data.

\vspace{0.25cm}
\begin{center}
\begin{tabular}{>{\RaggedRight\arraybackslash}p{0.28\linewidth} >{\centering\arraybackslash}p{0.26\linewidth} >{\centering\arraybackslash}p{0.26\linewidth}}
\toprule
Measure & Sample statistic & Population parameter \\
\midrule
Mean & $\bar{x}$ & $\mu$ \\
Proportion & $\hat{p}$ & $p$ \\
\bottomrule
\end{tabular}
\end{center}

\vspace{0.2cm}
The symbols matter because they tell us whether a number comes from the sample or from the full population.
\end{frame}

\begin{frame}{Population, Sample and Sampling Distribution}
Three related ideas appear throughout this chapter.
\begin{itemize}
  \item The \textbf{population distribution} describes values for individuals in the population
  \item The \textbf{sample distribution} describes the observations in one particular sample
  \item The \textbf{sampling distribution} describes the values of a statistic across many possible samples
\end{itemize}

\vspace{0.2cm}
The most common confusion in this chapter is mixing up the distribution of individual observations with the distribution of a statistic such as $\bar{x}$.
\end{frame}

\begin{frame}{Population, Sample and Sampling Distribution}
  \centering
  \lectureimagefull{fig_ch13_three_distributions.png}
\end{frame}

\begin{frame}{Example 13.1: A Population Parameter}
\begin{block}{Question}
A national registration system contains birth records for every birth in the population of interest during one year. The records show that 8\% of births were low birth weight and that the mean age at first birth was 26.3 years. Are these values parameters or statistics
\end{block}

\pause

\begin{block}{Explanation}
These values were calculated from the entire population defined for the study rather than from a sample.
\end{block}

\pause

\begin{block}{Solution}
Both values are \textbf{parameters}.
\[
p = 0.08
\]
\[
\mu = 26.3
\]
\end{block}
\end{frame}

\begin{frame}{Example 13.2: A Sample Statistic}
\begin{block}{Question}
A polling organization interviews a random sample of 1023 adults. In the sample, 19\% report smoking cigarettes during the previous week. Among the sampled smokers, the mean number of cigarettes per day is 12. Are these values parameters or statistics
\end{block}

\pause

\begin{block}{Explanation}
The values were computed from a sample rather than from every adult in the target population.
\end{block}

\pause

\begin{block}{Solution}
Both values are \textbf{statistics}.
\[
\hat{p} = 0.19
\]
\[
\bar{x} = 12
\]
\end{block}
\end{frame}

\section{The Sampling Distribution of \texorpdfstring{$\bar{x}$}{xbar}}

\begin{frame}{Repeated Sampling and the Sampling Distribution}
A sample mean is not fixed before the sample is selected. It depends on which individuals happen to be chosen.

\vspace{0.2cm}
A \textbf{sampling distribution} answers a repeated-sampling question.
\begin{block}{Sampling Distribution}
The sampling distribution of a statistic is the probability distribution of that statistic over all possible random samples of the same size drawn under the same sampling process.
\end{block}

A computer simulation can approximate this idea by drawing many samples and calculating the statistic each time.
\end{frame}

\begin{frame}{Examples 13.3 and 13.4: Simulating Sample Means}
\begin{block}{Question}
Individual IQ scores are modelled by a Normal distribution with mean $\mu = 100$ and standard deviation $\sigma = 15$. What happens if 500 random samples of size $n = 25$ are selected and the sample mean is calculated for each sample
\end{block}

\pause

\begin{block}{Explanation}
Each sample contributes one number to the sampling distribution, namely its sample mean.

Because the population is Normal, the sampling distribution of $\bar{x}$ is also Normal.
\[
\mu_{\bar{x}} = 100
\]
\[
\sigma_{\bar{x}} = \frac{15}{\sqrt{25}} = 3
\]
\end{block}
\end{frame}

\begin{frame}{Examples 13.3 and 13.4: Simulating Sample Means}
\begin{block}{Solution}
The histogram of the 500 simulated means is centred near 100 and has much less spread than the distribution of individual IQ scores.
\end{block}
\end{frame}

\begin{frame}{Examples 13.3 and 13.4: Simulated Sampling Distribution}
\centering
\lectureimagefull{fig_ch13_sampling_dist.png}
\end{frame}

\begin{frame}{The Sampling Distribution of $\bar{x}$}
If a random sample of size $n$ is drawn from a population with mean $\mu$ and standard deviation $\sigma$, then the sampling distribution of $\bar{x}$ has two key properties.

\vspace{0.2cm}
\textbf{Centre}
\[
\mu_{\bar{x}} = \mu
\]

\textbf{Spread}
\[
\sigma_{\bar{x}} = \frac{\sigma}{\sqrt{n}}
\]

\vspace{0.2cm}
The sample mean is therefore an \textbf{unbiased estimator} of the population mean. Larger samples make the sampling distribution narrower.
\end{frame}

\begin{frame}{The Effect of Sample Size}
  \centering
  \lectureimagefull{fig_ch13_sample_size.png}
\end{frame}

\begin{frame}{Why the Square Root Relationship Matters}
The standard deviation of the sample mean is
\[
\sigma_{\bar{x}} = \frac{\sigma}{\sqrt{n}}
\]

\pause

If the sample size is multiplied by 4, then
\[
\frac{\sigma}{\sqrt{4n}} = \frac{1}{2}\frac{\sigma}{\sqrt{n}}
\]

\pause

If the sample size is multiplied by 100, then
\[
\frac{\sigma}{\sqrt{100n}} = \frac{1}{10}\frac{\sigma}{\sqrt{n}}
\]

\pause

A larger sample improves precision, but the improvement is slower than many students expect.
\end{frame}

\section{Shape and the Central Limit Theorem}

\begin{frame}{When Is the Sampling Distribution Normal}
The shape of the sampling distribution of $\bar{x}$ depends on the population and the sample size.

\vspace{0.2cm}
Two important cases appear in first-year statistics.
\begin{itemize}
  \item If the population is Normal, then the sampling distribution of $\bar{x}$ is exactly Normal for any sample size
  \item If the population is not Normal, the Central Limit Theorem tells us that the sampling distribution of $\bar{x}$ becomes approximately Normal when the sample size is sufficiently large
\end{itemize}
\end{frame}

\begin{frame}{Example 13.5: Population Values and Sample Means}
\begin{block}{Question}
Individual IQ scores are modelled by a Normal distribution with $\mu = 100$ and $\sigma = 15$. Random samples of size $n = 25$ are selected. How does the population distribution compare with the sampling distribution of $\bar{x}$
\end{block}

\pause

\begin{block}{Explanation}
Both distributions are centred at 100. The difference is the spread.
\[
\sigma = 15
\]
\[
\sigma_{\bar{x}} = \frac{15}{\sqrt{25}} = 3
\]
\end{block}
\end{frame}

\begin{frame}{Example 13.5: Population Values and Sample Means}
\begin{block}{Solution}
The sampling distribution of sample means is much narrower than the population distribution of individual scores.
\end{block}
\end{frame}

\begin{frame}{Example 13.5: Population and Sampling Distributions}
\centering
\lectureimagefull{fig_ch13_pop_vs_samp.png}
\end{frame}

\begin{frame}{The Central Limit Theorem}
\begin{block}{Central Limit Theorem}
For independent observations from a population with finite mean $\mu$ and finite standard deviation $\sigma$, the sampling distribution of $\bar{x}$ becomes approximately Normal as the sample size becomes large.
\end{block}

\vspace{0.2cm}
The theorem does \textbf{not} say that the raw observations in one large sample become Normal.

\vspace{0.2cm}
The theorem does say that the distribution of \textbf{sample means across repeated samples} becomes more nearly Normal as $n$ increases.
\end{frame}

\begin{frame}{Example 13.6: The Central Limit Theorem in Action}
\begin{block}{Question}
A population is strongly right-skewed. How does the shape of the sampling distribution of $\bar{x}$ change as the sample size increases
\end{block}

\pause

\begin{block}{Explanation}
A sample mean averages several observations. As more observations are included, the distribution of the mean becomes less skewed.
\end{block}
\end{frame}

\begin{frame}{Example 13.6: The Central Limit Theorem in Action}
\begin{block}{Solution}
The population remains right-skewed. The sampling distribution of the sample mean becomes more symmetric and more nearly Normal as the sample size grows.
\end{block}
\end{frame}

\begin{frame}{Example 13.6: The Central Limit Theorem in Action}
\centering
\lectureimagefull{fig_ch13_clt_action.png}
\end{frame}

\begin{frame}{Assessing Whether a Normal Approximation Is Reasonable}
A practical first-year approach is to consider three things together.
\begin{itemize}
  \item The shape of the population or sample histogram
  \item The sample size
  \item Whether severe outliers or strong skewness are present
\end{itemize}

\vspace{0.2cm}
A mild departure from symmetry often requires less averaging than a severe departure from symmetry.
\end{frame}

\begin{frame}{Illustrative Shapes for CLT Discussion}
  \centering
  \lectureimagefull{fig_ch13_practical.png}
\end{frame}

\section{The Sampling Distribution of \texorpdfstring{$\hat{p}$}{phat}}

\begin{frame}{The Sample Proportion}
For a categorical variable with two outcomes, the statistic of interest is often the sample proportion.

\vspace{0.2cm}
If $X$ is the number of successes in a sample of size $n$, then
\[
\hat{p} = \frac{X}{n}
\]

\vspace{0.2cm}
The word success is only a label for the outcome being counted.
\end{frame}

\begin{frame}{The Sampling Distribution of $\hat{p}$}
The sampling distribution of $\hat{p}$ has properties similar to the sampling distribution of $\bar{x}$.

\vspace{0.2cm}
\textbf{Centre}
\[
\mu_{\hat{p}} = p
\]

\textbf{Spread}
\[
\sigma_{\hat{p}} = \sqrt{\frac{p(1-p)}{n}}
\]

\vspace{0.2cm}
A Normal approximation is typically used when both expected counts are at least 10.
\[
np \ge 10
\]
\[
n(1-p) \ge 10
\]
\end{frame}

\begin{frame}{Example 13.8: Contact Lens Practices}
\begin{block}{Question}
A random sample of 281 contact-lens users contains 139 people who do not consistently wash their hands before handling their lenses. What is the sample proportion
\end{block}

\pause

\begin{block}{Explanation}
The count of successes is $X = 139$ and the sample size is $n = 281$. The sample proportion is the count divided by the sample size.
\end{block}

\pause

\begin{block}{Solution}
\[
\hat{p} = \frac{139}{281}
\]
\[
\hat{p} \approx 0.495
\]
The sample proportion is about 49.5\%.
\end{block}
\end{frame}

\begin{frame}{Examples 13.9 and 13.10: Blood Type O}
\begin{block}{Question}
Suppose the population proportion with blood type O is $p = 0.45$. A random sample of $n = 1000$ records is selected. What is the probability that the sample proportion is 0.42 or lower
\end{block}

\pause

\begin{block}{Step 1: Identify the Sampling Distribution}
The centre is
\[
\mu_{\hat{p}} = 0.45
\]
The standard deviation is
\[
\sigma_{\hat{p}} = \sqrt{\frac{0.45(0.55)}{1000}} \approx 0.01573
\]
The success-failure condition is satisfied because $np = 450$ and $n(1-p) = 550$.
\end{block}
\end{frame}

\begin{frame}{Examples 13.9 and 13.10: Blood Type O}
\begin{block}{Step 2: Convert the Question to a Probability}
The probability of interest is
\[
P(\hat{p} \le 0.42)
\]
\end{block}

\pause

\begin{block}{Step 3: Standardize}
\[
z = \frac{0.42 - 0.45}{0.01573} \approx -1.91
\]
\end{block}

\pause

\begin{block}{Solution}
Using the Normal model gives
\[
P(\hat{p} \le 0.42) \approx 0.0283
\]
This is about 2.83\%.
\end{block}
\end{frame}

\begin{frame}[fragile]{Examples 13.9 and 13.10: R Code}
The same probability can be calculated in R with \texttt{pnorm()}.

\scriptsize
\begin{verbatim}
# Population proportion and sample size
p_type_o <- 0.45
n_type_o <- 1000

# Standard deviation of the sampling distribution of p-hat
sd_type_o <- sqrt(p_type_o * (1 - p_type_o) / n_type_o)

# Lower-tail probability P(p-hat <= 0.42)
prob_type_o <- pnorm(
  q = 0.42,
  mean = p_type_o,
  sd = sd_type_o
)

print(round(prob_type_o, 4))
\end{verbatim}
\normalsize
\end{frame}

\begin{frame}{Sampling Distribution of $\hat{p}$}
  \centering
  \lectureimagefull{fig_ch13_phat_dist.png}
\end{frame}

\section{The Law of Large Numbers}

\begin{frame}{The Law of Large Numbers}
\begin{block}{Law of Large Numbers}
For independent observations with a finite expected value, the sample mean converges to the population mean as the sample size increases. For independent Bernoulli observations, the sample proportion converges to the population probability.
\end{block}

\vspace{0.2cm}
This is a long-run idea. It does not say that every new observation brings the statistic closer to the parameter. It says that the overall tendency is toward the true population value.
\end{frame}

\begin{frame}{The Law of Large Numbers in Action}
  \centering
  \lectureimagefull{fig_ch13_lln.png}
\end{frame}

\begin{frame}{Why the Law of Large Numbers Matters}
The Law of Large Numbers explains why averages become more stable in large collections of observations.

\vspace{0.2cm}
Examples include
\begin{itemize}
  \item Long-run average claim costs in insurance portfolios
  \item Long-run average outcomes in repeated physical or biological measurements
  \item Stabilization of a running proportion or a running mean in simulation studies
\end{itemize}

In the life sciences, sample sizes are often limited by cost, time or ethics, which is why sampling distributions are so important.
\end{frame}

\section{Practice Examples}

\begin{frame}{Practice Example: Sampling Distribution of $\bar{x}$}
\begin{block}{Question}
A genetic strain of laboratory mice has population mean gestation period $\mu = 20$ days and population standard deviation $\sigma = 1.5$ days. A random sample contains $n = 15$ mice. What are the mean and standard deviation of the sampling distribution of $\bar{x}$
\end{block}

\pause

\begin{block}{Explanation}
The centre of the sampling distribution equals the population mean. The spread is the population standard deviation divided by the square root of the sample size.
\end{block}

\pause

\begin{block}{Solution}
\[
\mu_{\bar{x}} = 20
\]
\[
\sigma_{\bar{x}} = \frac{1.5}{\sqrt{15}} \approx 0.387
\]
\end{block}
\end{frame}

\begin{frame}{Practice Example: Normal Population Probability}
\begin{block}{Question}
Blood cholesterol levels in healthy adult Golden Retrievers are modelled as Normal with mean $\mu = 200$ mg/dL and standard deviation $\sigma = 15$ mg/dL. A random sample contains $n = 9$ dogs. What is the probability that the sample mean exceeds 210 mg/dL
\end{block}

\pause

\begin{block}{Step 1: Describe the Sampling Distribution}
Because the population is Normal, the sampling distribution of $\bar{x}$ is Normal.
\[
\mu_{\bar{x}} = 200
\]
\[
\sigma_{\bar{x}} = \frac{15}{\sqrt{9}} = 5
\]
\end{block}
\end{frame}

\begin{frame}{Practice Example: Normal Population Probability}
\begin{block}{Step 2: State the Probability}
\[
P(\bar{x} > 210)
\]
\end{block}

\pause

\begin{block}{Solution}
Using the Normal model gives
\[
P(\bar{x} > 210) \approx 0.0228
\]
This is about 2.28\%.
\end{block}
\end{frame}

\begin{frame}[fragile]{Practice Example: Normal Population Probability R Code}
\scriptsize
\begin{verbatim}
# Sampling standard deviation for x-bar
sampling_sd <- 15 / sqrt(9)

# Upper-tail probability P(x-bar > 210)
prob_cholesterol <- pnorm(
  q = 210,
  mean = 200,
  sd = sampling_sd,
  lower.tail = FALSE
)

print(round(prob_cholesterol, 4))
\end{verbatim}
\normalsize
\end{frame}

\begin{frame}{Practice Example: CLT Probability}
\begin{block}{Question}
A population of viral-load measurements is strongly right-skewed with mean $\mu = 1000$ particles/mL and standard deviation $\sigma = 300$ particles/mL. A random sample contains $n = 50$ patients. What is the approximate probability that the sample mean is below 900 particles/mL
\end{block}

\pause

\begin{block}{Explanation}
The population is not Normal, but the sample size is reasonably large. The Central Limit Theorem supports an approximate Normal model for $\bar{x}$.
\[
\mu_{\bar{x}} = 1000
\]
\[
\sigma_{\bar{x}} = \frac{300}{\sqrt{50}} \approx 42.43
\]
\end{block}
\end{frame}

\begin{frame}{Practice Example: CLT Probability}
\begin{block}{Solution}
The probability of interest is
\[
P(\bar{x} < 900)
\]
Using the approximate Normal model gives
\[
P(\bar{x} < 900) \approx 0.0092
\]
This is about 0.92\%.
\end{block}
\end{frame}

\begin{frame}[fragile]{Practice Example: CLT Probability R Code}
\scriptsize
\begin{verbatim}
# Approximate sampling standard deviation from the CLT
sampling_sd_viral <- 300 / sqrt(50)

# Lower-tail probability P(x-bar < 900)
prob_viral <- pnorm(
  q = 900,
  mean = 1000,
  sd = sampling_sd_viral
)

print(round(prob_viral, 4))
\end{verbatim}
\normalsize
\end{frame}

\begin{frame}{Practice Example: Sample Proportion Probability}
\small
\begin{block}{Question}
An antibiotic has cure probability $p = 0.80$. A random sample contains $n = 200$ infected patients. What is the approximate probability that fewer than 75\% of the patients are cured
\end{block}

\pause

\begin{block}{Step 1: Check Conditions and Describe the Sampling Distribution}
The Normal approximation is appropriate because
\[
np = 160
\]
\[
n(1-p) = 40
\]
The centre and spread are
\[
\mu_{\hat{p}} = 0.80
\]
\[
\sigma_{\hat{p}} = \sqrt{\frac{0.80(0.20)}{200}} \approx 0.0283
\]
\end{block}
\end{frame}

\begin{frame}{Practice Example: Sample Proportion Probability}
\begin{block}{Step 2: State the Probability}
\[
P(\hat{p} < 0.75)
\]
\end{block}

\pause

\begin{block}{Solution}
Using the Normal model gives
\[
P(\hat{p} < 0.75) \approx 0.0385
\]
This is about 3.85\%.
\end{block}
\end{frame}

\begin{frame}[fragile]{Practice Example: Sample Proportion Probability R Code}
\scriptsize
\begin{verbatim}
# Model parameters for the sample proportion
p_cure <- 0.80
n_cure <- 200
sd_cure <- sqrt(p_cure * (1 - p_cure) / n_cure)

# Lower-tail probability P(p-hat < 0.75)
prob_cure <- pnorm(
  q = 0.75,
  mean = p_cure,
  sd = sd_cure
)

print(round(prob_cure, 4))
\end{verbatim}
\normalsize
\end{frame}

\begin{frame}{Preview of Statistical Inference}
\begin{block}{Question}
A baseline model gives mean recovery time $\mu = 14$ days with standard deviation $\sigma = 3$ days. A random sample of $n = 100$ patients receiving a new treatment has mean recovery time $\bar{x} = 13.1$ days. If the baseline model were still true, how unusual would a sample mean of 13.1 days or lower be
\end{block}

\pause

\begin{block}{Explanation}
Under the baseline model, the sampling distribution of $\bar{x}$ has
\[
\mu_{\bar{x}} = 14
\]
\[
\sigma_{\bar{x}} = \frac{3}{\sqrt{100}} = 0.30
\]
\end{block}
\end{frame}

\begin{frame}{Preview of Statistical Inference}
\begin{block}{Solution}
The probability of interest is
\[
P(\bar{x} \le 13.1)
\]
Using the Normal model gives
\[
P(\bar{x} \le 13.1) \approx 0.0013
\]
\end{block}

\pause

A small probability means that the observed sample mean would be unusual if the baseline model were true. This repeated-sampling logic is the basis of formal hypothesis testing in later chapters.
\end{frame}

\begin{frame}[fragile]{Preview of Statistical Inference R Code}
\scriptsize
\begin{verbatim}
# Baseline model for the sampling distribution of x-bar
baseline_mean <- 14
population_sd <- 3
sample_size <- 100
sampling_sd_recovery <- population_sd / sqrt(sample_size)

# Lower-tail probability P(x-bar <= 13.1)
prob_recovery <- pnorm(
  q = 13.1,
  mean = baseline_mean,
  sd = sampling_sd_recovery
)

print(round(prob_recovery, 4))
\end{verbatim}
\normalsize
\end{frame}

\begin{frame}{R Functions Used in This Lecture}
Several R functions appear in the notes and examples.
\begin{itemize}
  \item \texttt{pnorm()} for Normal probabilities
  \item \texttt{sqrt()} for square roots in standard-deviation formulas
  \item \texttt{round()} for rounded numerical output
  \item \texttt{rnorm()} for simulated Normal data in the notes
  \item \texttt{rbinom()} for simulated binomial counts in the notes
\end{itemize}

\vspace{0.2cm}
The lecture figures in the notes also use \texttt{ggplot()} and related plotting layers to visualize sampling distributions.
\end{frame}

\begin{frame}{Chapter Summary}
\begin{itemize}
  \item A parameter describes a population and a statistic is calculated from a sample
  \item A sampling distribution describes how a statistic behaves across repeated random samples
  \item For a sample mean, the centre is $\mu$ and the spread is $\sigma/\sqrt{n}$
  \item For a sample proportion, the centre is $p$ and the spread is $\sqrt{p(1-p)/n}$
  \item The Central Limit Theorem explains why sample means become approximately Normal for large samples
  \item The Law of Large Numbers explains why sample statistics stabilize near the true population value in the long run
\end{itemize}
\end{frame}

\end{document}
